\documentclass[11pt]{article}
\usepackage[a4paper,margin=28mm]{geometry}
\usepackage{amsmath,amssymb}
\usepackage{hyperref}
\setlength{\emergencystretch}{2em}
\title{A unique hypergraph decomposition at every prime order}
\author{Ziang Ni (\texttt{ziangni-sys})}
\date{10 October 2026}
\begin{document}
\maketitle
\noindent\textbf{Disproof of conjecture 00000008437.}
The conjecture claims double-exponential numbers of decompositions
of a complete $t$-uniform hypergraph into specified isomorphism
classes when the number $v$ of vertices is prime. It imposes no
exclusion of single-edge block classes. That permitted class has
exactly one decomposition at every order.

Take $t=2$. The complete $2$-uniform hypergraph $K_v^{(2)}$ has vertex
set $V=\{0,\ldots,v-1\}$ and edge set
\[
 E=\{e\subseteq V:|e|=2\}.
\]
Specify the block isomorphism class $H$ consisting of a single
$2$-edge on its two incident vertices. Each copy is indexed by one
edge $e\in E$, and its edge set is $\{e\}$. All such blocks are
isomorphic to $H$; unused ambient vertices are not part of a block's
vertex support.

A standard unordered decomposition is a collection $\mathcal D$ of
edge-disjoint copies of $H$ covering all of $E$. Coverage forces the
block $\{e\}$ to occur for each $e\in E$: any block containing $e$
has only one edge and hence must equal $\{e\}$. Conversely every block
of $\mathcal D$ has that form. Therefore necessarily
\[
 \mathcal D=\big\{\{e\}:e\in E\big\}.
\]
This collection is indeed a decomposition: its singleton edge sets
are pairwise disjoint and their union is $E$. Thus the exact number
of decompositions is
\[
 N_H(v)=1\qquad\text{for every }v.
\]
In particular this holds for every prime $v$, including arbitrarily
large primes. A constant function cannot have double-exponential
growth along the primes. For example, for every positive $c$, the
lower bound $N_H(v)\geq\exp(\exp(cv))$ fails at every positive prime;
more generally $N_H$ does not even tend to infinity. This refutes the
first conjunct, so no assertion about automorphism fixed-point
spectra is needed.

\paragraph{Scope.}
This uses an admissible specified isomorphism class in the actual
complete hypergraph. A theorem restricted to nontrivial multi-edge
blocks would require a hypothesis absent from both source languages.
If decompositions were artificially ordered, the count would instead
be $\binom v2!$; even that is bounded by
$\exp(O(v^2\log v))$ and is not double-exponential in $v$. The usual
unordered convention gives the stronger exact count above.

\paragraph{Lean coverage.}
The project represents actual edges as two-element subsets of
$\mathrm{Fin}(v)$, blocks as finite edge sets, and decompositions as
unordered finite families satisfying the one-edge block condition,
coverage, and pairwise disjointness. It proves existence and uniqueness,
the exact cardinality one, and counterexamples at arbitrarily large
primes. Lean 4.19.0 and Mathlib are publicly pinned; theorem audits
use only standard logical axioms.
\end{document}
